% TOP_PROBLEM: {"id": 30006945, "title": "Erdős Matching Conjecture", "category_id": 2, "category": {"id": 2, "name": "combinatorics", "display_name": "Combinatorics", "description": "Counting problems, graph theory, discrete structures.", "slug": "combinatorics", "order_index": 2, "created_at": "2026-07-31T15:26:25.671Z"}, "status": "open"}
% FIRST_POSTED: 2026-09-25
% ENTRY_KIND: statement_only
% !TeX program = lualatex
% Definition and sources only; this file is not an LLM solution attempt.
\documentclass[11pt]{article}
\usepackage[margin=25mm]{geometry}
\usepackage{fontspec}
\setmainfont{Latin Modern Roman}
\usepackage{amsmath,amssymb,mathtools}
\usepackage{unicode-math}
\setmathfont{Latin Modern Math}
\usepackage{xurl,hyperref}
\hypersetup{colorlinks=true,urlcolor=blue}
\setcounter{secnumdepth}{0}
\setlength{\parindent}{0pt}
\setlength{\parskip}{5pt}
\setlength{\emergencystretch}{3em}
\begin{document}
\section*{\texorpdfstring{Erdős Matching Conjecture}{Erdős Matching Conjecture}}\label{30006945}
\subsection{Definitions and mathematical statement}
Let $n,k,s$ be positive integers with $n\ge ks$, and let $\mathcal F\subseteq\binom{[n]}k$. Its matching number is the maximum number of pairwise disjoint members. If this number is at most $s-1$, the conjecture asserts
\[
 |\mathcal F|\le
 \max\left\{\binom{ks-1}{k},
 \binom nk-\binom{n-s+1}{k}\right\}.
\]
The first candidate is realized by all $k$-subsets of a fixed $(ks-1)$-element set. The second is realized by all $k$-subsets meeting a fixed $(s-1)$-element set. Each family avoids $s$ disjoint members, explaining both terms and the parameter convention. Formulations using matching number at most $s$ replace this entry's $s$ by $s+1$; mixing those conventions changes the formula.
\par\smallskip\textit{Formulation and concept references:} \href{https://arxiv.org/pdf/2602.19230}{[S1]}.

\subsection{Short English statement}
How large can a uniform family of sets be when it contains no s pairwise disjoint members? The conjecture gives two competing sharp constructions.
\subsection{Sources}
\begin{itemize}
\item[S1]Frankl, Lu, Ma and Wu, Towards the Erdős matching conjecture for 4-uniform hypergraphs, Conjecture1.1; substitute q=s-1 in their matching-number convention.
\url{https://arxiv.org/pdf/2602.19230}.
\textit{Formulation source.}
\end{itemize}
\end{document}
